\begin{textsample}{NEG Dim 6 – human – Score -2.00 – sp/sp\_ef\_matematica\_2\_p7.txt}  \label{ex:f6_neg_013}
Raciocinar matematicamente oportuniza desenvolver algumas formas de pensar muito próprias da Matemática, dentre as quais destacam -se o pensar indutivo, o dedutivo, o espacial e o não determinístico.
Essas diferentes formas de pensar contribuem para que os estudantes aprendam a raciocinar a partir das evidências que encontram em suas explorações e investigações e do que já sabem que é verdade.
Aprendam, ainda, a reconhecer as características de uma ideia aceitável em Matemática, desenvolvendo raciocínios cada vez mais sofisticados, tais como análise, prova, avaliação, explicação, inferência, justificativa e generalização, dependendo da situação-problema que enfrentam.

% matched lemmas: 
\end{textsample}
